\chapter{Introduction}
\label{ch:intro}

\section{Motivation}

Bangla is one language with many voices. A speaker from Sylhet, one from
Chattogram and one from Barishal each speak recognisably differently from one
another and from the standard register heard on national news, in pronunciation,
in rhythm and intonation, and in vocabulary. To a listener from Bangladesh these
differences are often audible within a few seconds. For speech technology they
matter twice: dialect is a source of error for speech recognisers trained on
standard speech, and it is information in its own right, useful for routing
calls, adapting recognisers, or studying language variation at scale.

Automatic \emph{dialect identification} (DID) is the task of predicting a
speaker's dialect from their speech. It sits between language identification,
where the classes are different languages, and speaker identification, where
the classes are individuals \citep{snyder2018xvectors}. For widely studied
languages such as Arabic it has dedicated benchmarks \citep{ali2019mgb5}; for
Bangla it does not. Bangla speech processing is low-resource in general
\citep{bengaliasrsurvey2021}, and previous Bangla dialect work has relied on
hand-crafted acoustic features or on written dialect text
\citep{bengaliaccent2021,bangladial2025}.

Two developments make the task newly tractable. Self-supervised and weakly
supervised speech models---wav2vec~2.0, XLS-R and Whisper
\citep{baevski2020wav2vec,babu2022xlsr,radford2023whisper}---learn general
representations from very large amounts of multilingual audio, including Bangla,
and transfer to new tasks with little labelled data \citep{yang2021superb}. And
public Bangla speech resources have grown: a 1,178-hour corpus of standard
Bangla \citep{rakib2023oodspeech} and, recently, a 78-hour corpus of regional
speech with transcripts \citep{dipto2025ben10}.

\section{Problem statement}

This thesis classifies 5-second audio clips into nine classes: eight regional
dialects (Barishal, Chattogram, Dhaka, Khulna, Mymensingh, Noakhali, Rajshahi,
Sylhet) and the standard register, labelled \emph{Formal}. The corpus contains
9,303 clips, about 13.1 hours, collected from television, film and online media
(\cref{ch:dataset}).

Found data of this kind brings a specific risk. The corpus has no speaker labels,
and for most clips no record of which programme or recording they came from. A
random split can therefore place clips of the same speaker, recorded with the
same microphone in the same room, in both the training and the test set. A model
can then score well by recognising the speaker or the recording channel rather
than the dialect. This is a known failure mode of audio classification on found
data \citep{geirhos2020shortcut,echohawk2026}, but it is usually disclosed as a
caveat rather than measured. Measuring it is a central aim of this thesis.

\section{Two routes to dialect identity}

Dialect is carried by both sound and words. Some cues are acoustic: vowel
quality, consonant realisation, pitch contour and tempo. Others are lexical:
a Barishal speaker may say \bn{mui}{মুই} where standard Bangla has \bn{ami}{আমি} (``I''); a
Sylhet speaker asks \bn{kita}{কিতা} where the standard form is \bn{ki}{কী} (``what'').
The acoustic models used in most of this thesis see only the sound: they consume
log-mel spectrogram frames (Whisper) or the raw waveform (XLS-R) and have no
access to what was said.

The thesis therefore compares two routes:
\begin{itemize}
  \item the \textbf{acoustic route}, which learns dialect from the audio signal
        with hand-crafted features, frozen pretrained embeddings and an
        end-to-end fine-tuned model; and
  \item the \textbf{lexical route}, which transcribes each clip with
        off-the-shelf Bangla speech recognition, separates regional words from
        standard ones, and identifies the dialect from the words.
\end{itemize}
Combining the two answers the question that gives the thesis its title: how much
of Bangla dialect identity is carried by the words, and how much lies beyond
them?

\section{Research questions}

\begin{enumerate}[label=\textbf{RQ\arabic*.}, leftmargin=*]
  \item How accurately can pretrained speech models identify nine Bangla dialect
        classes from 5-second clips, compared with hand-crafted acoustic
        features?
  \item How much of that accuracy comes from speaker and recording-channel
        overlap between training and test data, and do the conclusions survive a
        speaker-disjoint evaluation?
  \item How much dialect information survives off-the-shelf Bangla speech
        recognition, and how well can the dialect be identified from the
        resulting words alone?
  \item Are the acoustic and lexical routes complementary?
\end{enumerate}

\section{Contributions}

\begin{enumerate}
  \item A 9-class, 13.1-hour Bangla dialect corpus pipeline with
        validation, standardisation and leakage-aware splitting, released with
        its exact train/validation/test splits (\cref{ch:dataset,ch:method}).
  \item An application of self-supervised and weakly supervised multilingual
        speech models (Whisper, XLS-R) to Bangla regional dialect
        identification, compared with a hand-crafted baseline under one
        evaluation protocol (\cref{sec:res-clip}).
  \item A measurement of speaker and channel leakage. A strict split built from
        ECAPA-TDNN pseudo-speaker clusters, with its clustering threshold
        calibrated against the known sources of the Formal clips, turns the
        usual ``upper bound'' disclaimer into a measured effect of 12--14
        macro-F1 points for the pretrained models (\cref{sec:res-strict}).
  \item An ablation of representation against split protocol, speaker-group
        bootstrap confidence intervals for every score and gap, and a
        data-efficiency analysis (\cref{sec:res-strict,sec:res-efficiency}).
  \item A negative result: channel-suppression augmentation lowers
        strict-split accuracy, which the thesis attributes to prosody being
        dialect signal rather than nuisance (\cref{sec:res-augment,sec:disc-augment}).
  \item The lexical route: five off-the-shelf Bangla speech recognisers compared
        by how much dialect they preserve, using an automated dialect-word
        recall over 1,666 human-transcribed clips; regional-word matching with
        an interpretable per-region match percentage; a text classifier; and
        acoustic--lexical fusion, which raises the strict-split score from
        0.808 to 0.856 (\cref{sec:res-lexical,sec:disc-words}).
\end{enumerate}

No speech-recognition model was trained or fine-tuned in this work. The only
model trained end to end is the XLS-R dialect classifier; the speech
recognisers are used exactly as their publishers released them.

\section{Thesis outline}

\Cref{ch:background} reviews dialect identification, pretrained speech
representations, Bangla speech resources and evaluation leakage.
\Cref{ch:dataset} describes the corpus. \Cref{ch:method} sets out the splitting
protocols, the acoustic models and the lexical route. \Cref{ch:results} reports
the results, \cref{ch:discussion} interprets them and states the limitations,
and \cref{ch:conclusion} answers the research questions.
\Cref{app:repro,app:settings} give the commands and settings needed to
reproduce every number.
